% =====================================================================
%  EXPERIMENTAL SETUP  —  standalone working draft (v4, post 3 review rounds)
%  Fresh draft (companion to / eventual replacement for methodology.tex).
%
%  Dual purpose:
%    (1) preliminary version of the paper's Experimental Setup section;
%    (2) an internal GUIDE — for each experiment we separate what we
%        want to COLLECT (the full arm matrix) from what we intend to
%        SHOW in the paper (a curated subset) and say WHY.
%
%  Convention (from methodology.tex): NO result numbers here; "What it
%  shows" states the intended claim / diagnostic, not a value.
%
%  \sysname{} = the framework name. The repo is inconsistent
%  (prose "McBenchface", code/tables "Choreo"); define ONCE upstream.
\providecommand{\sysname}{Choreo}
%  Requires: siunitx, booktabs, wasysym (for \CIRCLE \LEFTcircle \Circle).
%
%  RECONCILE (cross-file, HIGH): methodology.tex §exp-topology still says
%  "We do NOT run a separate independent-workload interference experiment
%  ... left to future work" and carries reviewer TODO D1. This draft now
%  DELIVERS that experiment (E6) and the MLPerf contrast (E5). When this
%  file replaces the methodology's Experiments subsection, retire that
%  paragraph and soften the Intro's scalability wording to match E7.
% =====================================================================

\section{Experimental Setup}
\label{sec:setup}

We evaluate \sysname{} with seven experiments in three tiers.
Two \emph{instrument-fidelity} experiments bound the measurement overhead
the framework itself introduces, so that any overhead we later attribute
to a workload is the workload's and not the harness's
(E1--E2). Two \emph{application case studies} then use \sysname{} to
surface end-to-end behaviour that single-stream, single-engine benchmarks
cannot observe --- a unified-memory bandwidth-contention study (E3) and a
pipeline-topology study of agentic RAG (E4). Finally, three
\emph{capability} experiments substantiate the framework's distinguishing
claims directly: that the prevailing methodology (MLPerf Inference) is a
\emph{special case} of \sysname{}'s configuration space (E5), that
collocated concurrent pipelines are a \emph{first-class, measurable}
regime (E6), and that stressing hardware of any capacity is a
\emph{no-code} change (E7). E5--E7 are the direct response to the concern
that collocation-aware profiling and no-code scalability, headlined in the
introduction, must be \emph{demonstrated} rather than asserted.

Because this document also guides our own data collection, each experiment
adds to its \textbf{Question}, \textbf{Design}, \textbf{Configurations},
and \textbf{Metrics} a \textbf{Collect vs.\ report} note separating the
full arm matrix we run from the curated subset we intend to publish, with
the rationale for the cut. Results are deferred to
Section~\ref{sec:results}; nothing below states a measured value.

% ---------------------------------------------------------------------
\subsection{Devices Under Test}
\label{sec:setup-duts}

We use two single-node machines from opposite extremes of the
unified-memory design space, chosen to test whether a measured systems
effect is intrinsic to a workload or specific to one platform, and to
exercise \emph{portability} across the ``mid-range gap'' between the rigid
\emph{Tiny}/\emph{Edge}/\emph{Datacenter} tracks of existing suites
(Section~\ref{sec:background}). We do \emph{not} treat the pair as a
controlled capacity axis: capacity differs by \num{8}$\times$ but
bandwidth by only ${\approx}\num{1.4}\times$, and device class, runtime,
and numerical precision all co-vary.

\begin{description}
  \item[Apple M2~Pro (edge / consumer).] A \SI{16}{\giga\byte}
    unified-memory laptop SoC in which GPU, Neural Engine, and CPU share
    one physical memory and its bandwidth
    ([[FILL: core split; GPU cores; ${\approx}\SI{200}{\giga\byte\per\second}$]]).
    Three backends over that shared memory: GPU via Apple MLX/Metal
    (PyTorch \texttt{mps}), Neural Engine (\texttt{ane}) via Core~ML, and
    CPU. Software: Python~3.10, \texttt{mlx}/\texttt{mlx-lm}, Core~ML
    Tools, \texttt{faiss-cpu} ([[FILL: macOS + package versions]]). LLM
    weights are 4-bit quantized to fit the \SI{16}{\giga\byte} budget
    alongside other resident stages.
  \item[NVIDIA DGX~Spark (GB10, large unified memory).] A
    Grace--Blackwell platform with \SI{128}{\giga\byte} LPDDR5X
    ([[FILL: ${\approx}\SI{273}{\giga\byte\per\second}$]]), presenting one
    CUDA GPU. Grace CPU and Blackwell GPU are separate dies joined by a
    coherent NVLink-C2C interconnect --- a coherent unified-memory
    \emph{system} rather than one physical pool. Software: PyTorch
    \texttt{2.10.0+cu130}, \texttt{transformers~5.2}, BF16.
\end{description}

\sysname{} runs the same declarative pipeline definitions on both
machines, changing only the per-stage \texttt{device}/backend field; no
code changes are required (Section~\ref{sec:framework}). Because the
M2~Pro runs 4-bit MLX weights and DGX~Spark runs BF16 PyTorch weights, the
two are not numerically identical models: we draw cross-device conclusions
only about the \emph{direction} of a systems effect, never about absolute
latency or quality across devices. All weights and datasets are
pre-fetched before measured runs. \textbf{Reproducibility.} Models are
pinned by Hugging Face commit revision and datasets by version/snapshot
date; the 4-bit MLX weights are produced from those revisions with a fixed
quantization recipe ([[FILL: Qwen3.5-9B/4B + CLIP-ViT-L/14 hashes; MLX
quant recipe; dataset snapshots; FAISS/ChromaDB index params]]). We
monitor package power/temperature (\texttt{macmon};
\texttt{nvidia-smi}/\texttt{dcgmi}, sampled at [[FILL: period]]) and
\emph{exclude} any run whose sustained clock or package temperature
crosses [[FILL: pre-registered throttle threshold]]; the threshold and
discard-vs-report policy are fixed before collection to remove
researcher-degrees-of-freedom.

% ---------------------------------------------------------------------
\subsection{Instrumentation and Metrics}
\label{sec:setup-metrics}

\paragraph{Measurement pipeline.} Every run emits a \emph{timing trace}
(each stage and the driver log \texttt{start}/\texttt{end} for their
\texttt{prepare}/\texttt{run} phases; per-query events carry the query id,
submission time, epoch, and batch index) and, for the case studies, a
per-query JSONL sidecar written by \texttt{TerminalCapture} (question,
golden answers, retrieved documents, generated answer). All durations use
a single monotonic clock (\texttt{perf\_counter\_ns};
resolution \SI{41.7}{\nano\second} on M2~Pro, \SI{1}{\nano\second} on
GB10); wall-clock time is used only for cross-process trace alignment.
RadT attaches per-\emph{process} resource listeners (\texttt{macmon} on
Apple; \texttt{nvidia-smi}, \texttt{dcgmi}, \texttt{iostat}, \texttt{top}
on NVIDIA), sampled at [[FILL: period]]; because isolation is at the
process boundary (one process per pipeline, stages as threads within it),
these attribute resource use per pipeline even under collocation. All
metrics are computed offline by \texttt{evaluation/scripts/}.

\paragraph{Performance metrics.} \textbf{Per-query latency} (run-start to
run-end); \textbf{throughput} (completed queries over makespan --- MLPerf's
\emph{Offline} throughput, not SLO-constrained \emph{Server});
\textbf{stage concurrency} (sweep-line count of overlapping \texttt{run}
intervals --- temporal overlap, not CPU parallelism, since pure-Python work
is GIL-serialized while native kernels overlap); \textbf{per-device
occupancy} (union of the run intervals of stages mapped to each device, as
a fraction of makespan --- residency, not hardware utilization; can exceed
100\% in aggregate under collocation). For the topology and collocation
studies we additionally report \textbf{LLM calls per query} (count of
generation invocations a query triggers, including retry/rewrite loops)
and \textbf{resident model memory} (peak device memory attributable to
loaded model weights + KV cache, from the resource listeners), the latter
being what the ${\approx}3\times$-weights cost of physical decomposition is
measured against.

\paragraph{Framework-overhead metrics.} Per-stage overhead (per-query
latency / depth, \si{\micro\second}), inter-stage transition cost (gap
between one stage's \texttt{end} and the next stage's \texttt{start}),
and, for the modularity study, relative overhead
$(t_{\text{\sysname}}-t_{\text{base}})/t_{\text{base}}$ against a
hand-written baseline.

\paragraph{Quality metrics (RAG/VQA).} The standard accuracy metric per
task family: \textbf{exact match (EM)} and \textbf{token-level F1} for
text RAG, and \textbf{VQA accuracy} (soft match over the annotator set)
for OK-VQA. These scorers are the gate on every accuracy-conditioned
comparison below and are \emph{required before those tables ship}.
% BLOCKER(collect): EM/F1 + VQA-accuracy scorers are NOT yet implemented;
% the released analysis computes only the containment proxy. No
% accuracy-conditioned result (E3, E4) may be reported until they exist.
As a coarse secondary signal we also compute an \textbf{answer-containment
rate} (normalized golden answer is a substring of the normalized
generation); because containment is gameable by verbosity and the topology
arms differ in verbosity, we report it only alongside mean answer length
and never as the accuracy of record. We also report \textbf{answered rate}
(non-empty, non-retry-exhausted) and a \textbf{cross-pipeline parity}
score (mean token-level Jaccard on shared questions; count at
Jaccard${\geq}0.5$) as an agreement --- not correctness --- check.

\paragraph{Statistics.} Each configuration is run $R{=}5$ times with
fixed, distinct seeds controlling query order and the Poisson arrival
process. \emph{The unit of replication is the run, not the query.} Within-run
per-query latencies are correlated (shared warm state, thermal drift,
scheduler phase), so we do \emph{not} pool raw queries and bootstrap them
as if independent --- that pseudoreplication understates variance. Instead
we summarise each run by its statistic of interest (e.g.\ median per-query
latency, or median per-step latency in E2), report the across-run
distribution over the $R{=}5$ run-level values, and form CIs with a
\emph{hierarchical (cluster) bootstrap} that resamples runs first and then
queries within a resampled run ($10^4$ resamples); we report run-to-run
variance explicitly and print the $R{=}5$ raw run-level values alongside
every CI, since with only five clusters the interval alone is not to be
over-trusted and the actual spread should be visible. Because $R{=}5$ is small, we treat run-level CIs as
indicative and, where a paired design exists (E2's interleaved arms,
shared seeds), also report a paired across-run difference, which is more
powerful for small fixed effects than comparing two marginal intervals.
For rate metrics (answered, containment, EM/accuracy) we report Wilson
95\% intervals on the pooled counts \emph{and} a run-level interval over
the $R{=}5$ per-run rates; because pooled Wilson ignores run-to-run rate
variance it is anti-conservative when that variance is non-trivial, and
the run-level interval then governs. Tail quantiles (p95) are reported
only for cells whose pooled sample reaches $\geq\num{500}$ queries; p99 is
not estimated at our sample sizes. Generation is greedy (\texttt{do\_sample=false},
temperature~0) for every arm of a comparison, verified with a repeat-run
identical-output check, so sampling is not a confound.

\paragraph{Experiment overview.} Table~\ref{tab:exp-overview} lists the
seven experiments, their tier, device(s), and the single claim each is
designed to support.

\begin{table}[t]\centering\small
\caption{The seven experiments. Tier: I = instrument-fidelity,
A = application case study, C = capability demonstration.}
\label{tab:exp-overview}
\begin{tabular}{@{}llll@{}}
\toprule
Exp & Tier & Device(s) & Claim it supports \\
\midrule
E1 NoOp overhead      & I & M2~Pro          & harness cost is small, flat, $O(1)$-copy \\
E2 Modularity         & I & M2~Pro, GB10    & wrapping real work is $\ll$ the work \\
E3 VQA contention     & A & M2~Pro          & heterogeneity gain is contention-bound \\
E4 Self-RAG topology  & A & M2~Pro, GB10    & decomposition is a real, config-only trade-off \\
E5 MLPerf reduction   & C & M2~Pro, GB10    & MLPerf scenarios are \sysname{} configurations \\
E6 Collocation scaling& C & M2~Pro, GB10    & inter-pipeline interference is first-class \& measurable \\
E7 Capacity sweep     & C & M2~Pro, GB10    & stressing any-size hardware is a no-code change \\
\bottomrule
\end{tabular}
\end{table}

% =====================================================================
\subsection{E1 --- Framework overhead (NoOp microbenchmark)}
\label{sec:e1-noop}

\textbf{Question.} \sysname{} runs every stage in its own thread and moves
data between stages through queues. Before trusting any end-to-end number
we bound that machinery: (i) how much latency does each added stage
contribute, and does it \emph{accumulate} with depth? (ii) does passing
data by reference avoid per-hop copy cost? (iii) what does the optional
MLflow tracing layer add on top of core dispatch?

\textbf{Design.} We build synthetic pipelines of pass-through (no-op)
stages that perform no computation, so all measured time is framework
overhead. Two sweeps, generated programmatically
(\texttt{noop\_chain\_generator.py}):
\begin{itemize}
  \item \emph{Depth sweep:} chain length
    $N\in\{1,\dots,10,16,32,50,64,100\}$, payload~0, isolating per-stage
    dispatch and queue hand-off. \emph{Crossed with a tracing factor}
    (off / on), giving the two arms below.
  \item \emph{Payload sweep:} fixed depth~10, payloads
    $\{0, \SI{1}{\kibi\byte}, \SI{1}{\mebi\byte}, \SI{10}{\mebi\byte}\}$
    under \emph{reference} vs.\ \emph{deep-copy} passing, isolating data
    movement. This sweep is run tracing-\emph{off} only (the copy cost is
    independent of the span layer).
\end{itemize}
The load generator is \emph{closed-loop}
(\texttt{OfflineLoadScheduler}, exactly one query in flight), so this
measures dispatch and reference-vs-copy cost in isolation and deliberately
not queueing (exercised in E3/E4/E6). The two tracing arms are: \emph{off}
(\texttt{CHOREO\_DISABLE\_TRACING=1}: core dispatch = thread wake + queue
hand-off + CSV log) and \emph{on} (three MLflow spans per stage per query).
$R{=}5$, \texttt{max\_queries}=101; we drop the first [[FILL: $k$; justify
$k$ absorbs import / thread-pool spin-up --- likely $>1$]] queries as
warm-up. Run on the M2~Pro only: with no accelerator kernel in a no-op
stage this is a property of the framework and CPython interpreter, not the
device, so a single quiet machine bounds it.
\emph{Provenance note (blocker):} the committed 185-CSV matrix was
collected on a busy shared aarch64 Linux box (visible as a depth-50
contention outlier), which contradicts the ``single quiet machine''
premise; it must be \textbf{re-collected on an idle M2~Pro} before any
number is quoted.

\textbf{Configurations.} \texttt{noop\_depth\_*\_mode\_ref.yml} $\times$
2 tracing arms (depth sweep) and
\texttt{noop\_d10\_s\{0,1024,1048576,10485760\}\_m\{ref,copy\}} tracing-off
(payload sweep); $R{=}5$. Committed matrix: 185 CSVs (110 tracing-off,
75 tracing-on --- the tracing-on set covers the 15-point depth sweep, not
the payload/copy cells).

\textbf{Metrics.} Median per-query latency, per-stage overhead
(\si{\micro\second}), inter-stage transition cost, reference-vs-copy
latency ratio. Read against three claims: \emph{(i) depth-flatness} ---
latency linear in depth, transition cost flat in depth; \emph{(ii)
zero-copy} --- reference passing $O(1)$, copy arm $O(\text{payload})$;
\emph{(iii) overhead-in-context} --- core dispatch, and the tracing layer,
are a negligible fraction of one real model or retrieval stage.

\textbf{Threats to validity.} No-op stages never release the GIL, so the
per-stage figure is a GIL-serialized \emph{worst case} and an upper bound
on native-kernel-bound stages; the copy arm measures in-process
\texttt{deepcopy} and therefore \emph{lower}-bounds process-level
serialization. The tracing-on cost on the \texttt{-p 0} path uses
\emph{async} local-store export, an async proxy for the case studies'
\emph{synchronous} server export; since claim (iii) licenses attributing
case-study overhead to the workload, we quote the canonical synchronous
tracing cost from an orchestrated run rather than the proxy
([[FILL: sync tracing cost/stage]]).

\textbf{Collect vs.\ report.} \emph{Collect:} depth$\times$tracing plus the
tracing-off payload$\times$copy sweep, re-collected on an idle M2~Pro.
\emph{Report:} the depth-flatness curve (transition cost vs.\ depth), the
zero-copy ref-vs-copy figure, and one overhead-in-context line using the
\emph{synchronous} tracing cost against a real LLM-stage latency. No device
sweep (instrument property, not portability).

% =====================================================================
\subsection{E2 --- Modularity overhead (real workload vs.\ monolith)}
\label{sec:e2-modularity}

\textbf{Question.} Empty stages being cheap (E1) does not prove that
wrapping a \emph{real} workload in \sysname{}'s graph/queue/thread
structure is cheap. Does expressing a training workload as a \sysname{}
pipeline inflate per-step latency relative to a hand-written monolith?

\textbf{Design.} One workload, two implementations: (a) a standalone
PyTorch script (\texttt{baseline\_finetune.py}); (b) an equivalent
\sysname{} pipeline (\texttt{torchvision\_training.yml}). Workload:
transfer-learning fine-tuning of EfficientNetV2-S~\cite{tan2021efficientnetv2}
on Imagenette~\cite{imagenette} --- frozen backbone, replaced 10-class head,
Adam (lr $10^{-3}$), cross-entropy, batch~8, the same \num{12810}
trainable parameters. Both share identical data loading, run the
\texttt{train} split only, and call \texttt{torch.\{mps,cuda\}.synchronize()}
at step end so logged time reflects hardware completion; the measured
per-step bracket excludes data loading. To isolate the wrapper we set
\texttt{num\_workers}=0 in both arms (no prefetch overlap); the \sysname{}
data loader remains a separate, independently measured stage. Each run is
one continuous epoch (1{,}100 steps); we drop the first \num{200} as
warm-up, summarise each run by its median per-step latency, and
\emph{interleave} the arms run-by-run so both see identical conditions. As
in E1 we measure core dispatch (tracing off) and full instrumented cost
(tracing on); the baseline carries no tracing. Run on both DUTs, reading
overhead within a device.

\textbf{Configurations.}
\texttt{run\_modularity.py --device \{mps,cuda\} --runs 5 --max-batches 1100}
(\texttt{--num-workers 0}, interleaved, both tracing arms); baseline via
\texttt{baseline\_finetune.py --no-radt}. Tracing-on uses the
\texttt{async\_tracing} RadT branch (commit \texttt{3ba61cb}), whose
off-critical-path span export is what makes the tracing number realistic.

\textbf{Metrics.} Run-level median per-step latency per implementation
(across-run distribution, $n{=}5$), and \sysname{}'s overhead over the
baseline as both an \emph{absolute} cost (\si{\micro\second}) and a ratio,
reported as a \emph{paired} across-run difference (interleaved arms, shared
seeds) with its bootstrap CI. Reporting the absolute cost makes the
magnitude explicit: an overhead whose CI excludes zero is real yet, at
\si{\micro\second} scale against a multi-\si{\milli\second} step,
practically negligible.

\textbf{Threats to validity.} A modularity wrapper can only add work, so
overhead must be non-negative; a measured \emph{negative} value signals a
confound, not a speedup, and we treat it as a bug. We hit and fixed two:
(i) an earlier wall-clock reading differenced non-monotonic timestamps
across threads (fixed by the monotonic clock); (ii) an earlier schedule
(all baselines first, \texttt{num\_workers}=2 prefetch) let a data-path
difference dominate (fixed by interleaving and matching
\texttt{num\_workers}). Both are documented as evidence the metric is
correctly signed. Single batch size / model is deliberate ---
step-size-independence of the overhead is carried by E1.

\textbf{Collect vs.\ report.} \emph{Collect:} baseline + two tracing arms
$\times R{=}5$ on both DUTs (mps done; the \texttt{cuda}/GB10 half is
authoritative and [[FILL: confirm GB10 collection]]). \emph{Report:} the
per-step median table (baseline / off / on) with paired absolute + ratio
overhead per device. Until the GB10 half exists the portability read is
one-sided --- a blocker, not a footnote.

% =====================================================================
\subsection{E3 --- Unified-memory bandwidth contention (multimodal VQA)}
\label{sec:e3-vqa}

\textbf{Question.} On an Apple-silicon unified-memory SoC the GPU, Neural
Engine, and CPU share one memory and its bandwidth; single-engine
accelerator benchmarks cannot see the resulting contention. When we map
the stages of one workload across different engines, is the apparent
benefit of that heterogeneity real, or does it evaporate once queries
overlap and the engines contend for shared bandwidth? (M2~Pro only --- the
Neural Engine has no GB10 analogue.)

\textbf{Design.} A single image-grounded VQA pipeline over
OK-VQA~\cite{marino2019okvqa}: CLIP-ViT-L/14~\cite{radford2021clip} vision
encoding $\rightarrow$ FAISS~\cite{johnson2019faiss} retrieval (top-3) over
\num{10000} COCO-Karpathy captions~\cite{lin2014coco,karpathy2015deepvs}
$\rightarrow$ 4-bit MLX Qwen3.5-9B answer generation on the GPU. A
$2\times2$ over \emph{mapping} and \emph{schedule}:
\begin{itemize}
  \item \emph{Mapping~A (collocated):} CLIP and LLM both on GPU
    (\texttt{mps}); retrieval on CPU (\texttt{..._mapping\_a.yml}).
  \item \emph{Mapping~B (heterogeneous):} CLIP on the Neural Engine via
    Core~ML (\texttt{ane}); LLM on GPU; retrieval on CPU
    (\texttt{..._mapping\_b.yml}).
  \item \emph{Schedule:} \emph{pipelined} (Poisson \texttt{rate}=2.0~q/s,
    above the ${\approx}3$~s/query LLM service time) vs.\ \emph{serial}
    (\texttt{-{}-serialize}). \texttt{max\_queries}=30 [[FILL: raise toward
    the p95 gate ($\geq\num{500}$ pooled) or drop p95 for E3]].
\end{itemize}

\textbf{Configurations.} \texttt{multimodal\_vqa\_mapping\_\{a,b\}.yml}
$\times$ \{pipelined, serial\} $\times R{=}5$. We record measured
stage-concurrency mean per cell to confirm the entry queue actually backs
up (otherwise throughput at fixed Poisson rate reflects offered load, not
capacity).

\emph{Confound.} Mapping~B changes more than placement: the encoder
switches from PyTorch/Metal to a Core~ML FP16 package on the ANE, so
engine, runtime, and precision co-vary and the ANE has a first-call
compilation cost. We read the signal primarily \emph{within} a mapping
(serial vs.\ pipelined for the same encoder), discard the ANE first-call
query (reported separately), and use answer parity as a coarse divergence
check on the FP16 embeddings.

\textbf{Metrics.} The full performance set (\S\ref{sec:setup-metrics}) per
cell; VQA accuracy and answer parity for comparability. The diagnostic is
whether the heterogeneity advantage \emph{collapses} from serial to
pipelined. \emph{Mechanism attribution (strengthened):} occupancy is
residency, not utilization, and the ANE exposes no per-process utilization
counter, so occupancy/concurrency traces can show \emph{that} the
advantage collapses but cannot alone distinguish bandwidth contention from
command-buffer serialization. We therefore make a direct DRAM-bandwidth /
power read (\texttt{powermetrics}/\texttt{macmon}) a \textbf{required}
part of this experiment; the ``bandwidth-bound'' claim rests on it. Absent
that read we downgrade the claim to ``the heterogeneity gain collapses
under overlap, consistent with a shared-memory-contention regime.''

\textbf{Collect vs.\ report.} \emph{Collect:} the full $2\times2 \times
R{=}5$ plus the DRAM-bandwidth trace. \emph{Report:} the within-mapping
serial$\to$pipelined deltas for both mappings, the DRAM read that
adjudicates the mechanism, and accuracy+parity showing the arms answer the
same task. The raw B-vs-A delta appears only with the co-varying-factor
caveat.

% =====================================================================
\subsection{E4 --- Pipeline topology: monolithic vs.\ decomposed agentic RAG}
\label{sec:e4-selfrag}

\textbf{Question.} For an agentic, self-reflective RAG
workload~\cite{asai2024selfrag}, is one large model that performs every
sub-task (grading, answering, hallucination checking) in a single call
better than several smaller specialists composed into a graph? And ---
since ``decomposition'' bundles several changes --- how much of any effect
is \emph{model size}, how much \emph{physical} (separate resident
instances) vs.\ \emph{logical} (one shared instance) decomposition, and how
much the \emph{inference engine}? Such ablation-style restructuring is what
fixed-topology benchmarks cannot express but \sysname{}'s declarative,
non-linear graphs support directly.

\textbf{Design.} One Self-RAG task in a family of topologies that swap by
changing only configuration:
\begin{itemize}
  \item \emph{Monolith-9B:} one Qwen3.5-9B does grade+answer+check in a
    single structured pass; a router validates and optionally loops
    through a query rewriter (a bounded back-edge).
  \item \emph{Decomposed (3$\times$4B):} three separate Qwen3.5-4B
    instances (grader / generator / hallucination-checker; rewriter shares
    the grader) split the job, so cheap grade/check calls for one query
    overlap the expensive generation for another.
  \item \emph{Monolith-4B (size control):} one Qwen3.5-4B issues the three
    calls sequentially --- size-matched to Decomposed, isolating
    overlap/specialization at fixed parameter size.
  \item \emph{Decomposed-Shared (logical):} the same decomposed graph and
    per-role prompts, but all stages share \emph{one} resident 4B instance
    behind one mutex --- contrasted with Decomposed, isolates what the
    extra physical instances buy (cross-query overlap) against their
    ${\approx}3\times$ resident-memory cost.
\end{itemize}
The decomposed topology exercises fan-out routing, fan-in merge, and a
feedback cycle (retrieve--grade--rewrite), and is a within-pipeline
collocation case: multiple model instances co-resident on one accelerator,
profiled per stage. We run on two task difficulties --- factoid
(\texttt{rag-mini-wikipedia}, top-3 over 3{,}200 passages) and multi-hop
(HotpotQA~\cite{yang2018hotpotqa}, top-5 over a 20{,}000-passage corpus
from the gold+distractor contexts of the first 300 HotpotQA-dev
questions) --- retrieving from ChromaDB~\cite{chroma}, on both DUTs (GB10
BF16; M2~Pro 4-bit MLX). Retrieval differs across difficulties, so the
topology contrast is read \emph{within} a difficulty.

\emph{Engine comparison.} The two single-instance arms (Monolith-4B and
Decomposed-Shared) are additionally served on \emph{identical 4B weights}
by three backends --- HF Transformers (in-process), a local vLLM server
(continuous batching, via litellm), and a local Ollama server (via
litellm) --- to compare engines under identical work. vLLM is CUDA-only;
Ollama runs on both platforms.

\emph{Confounds and controls.} Monolith-9B-vs-Decomposed conflates size,
specialization, and call count; Monolith-4B (size) and Decomposed-Shared
(physical-vs-logical) are the controls that separate them.
\textbf{Because a multi-hop quality gap between 9B and 3$\times$4B is
exactly what a size difference alone would produce, the Monolith-4B
control must be run \emph{on multi-hop}, not only on factoid} --- otherwise
the headline ``quality-on-multi-hop'' result is unfalsifiable as a size
artifact. The four multi-hop control configs
(\texttt{multihop\_\{monolith\_4b,decomposed\_shared\}\_\{cuda,mlx\}}) do
not yet exist and are an explicit collection item.
% BLOCKER(collect): create multihop_monolith_4b_{cuda,mlx}.yml and
% multihop_decomposed_shared_{cuda,mlx}.yml — absent from configs/ today.

\emph{Load matching.} Load generation is not identical across DUTs: the
CUDA host uses \texttt{rate}=4.0 / 30 queries, the slower MLX host
\texttt{rate}=1.0 / 10 queries, each chosen so the entry queue backs up in
pipelined mode (confirmed by the measured stage-concurrency mean). On the
16\,GB M2~Pro the decomposed-pipelined cells can hit
\texttt{[METAL] Insufficient Memory} (3$\times$4B $\times$ in-flight KV
caches). \textbf{The fix must not confound topology with offered load:} we
re-run \emph{all} MLX topology arms at the same reduced load
(\texttt{rate}=0.4 / \texttt{max\_queries}=8) so every arm sees identical
arrivals, and we verify the queue still backs up at that rate; where it
does not, we restrict the MLX comparison to the serial cells and state the
pooled $N$ per arm (10--40 queries $\times R{=}5$, below the p95 gate --- so
no MLX tail latencies are reported). We expect this to be the common case:
at \texttt{rate}=0.4 against a seconds-scale LLM the entry queue usually
will not back up, so on the MLX host the \emph{pipelined} topology story
largely reduces to serial-only. We state this up front rather than as a
contingency --- the pipelined-overlap claim is carried by the GB10 cells.
Serial cells use \texttt{-{}-serialize}.

\textbf{Configurations.} The arm$\times$task$\times$backend matrix is in
\texttt{evaluation/self\_rag/configs/} (16 configs today; +4 multi-hop
controls to add). Each timing arm runs pipelined and serial $\times R{=}5$
per DUT. Table~\ref{tab:selfrag-matrix} lists arms and the report subset.

\begin{table}[t]\centering\small
\caption{Self-RAG arms. \CIRCLE\ = core (report); \LEFTcircle\ = extended
(collect, report if space/signal); \Circle\ = engine-comparison overlay;
\texttt{+} = config to be created. Every reported timing arm runs
pipelined + serial, $R{=}5$, per DUT.}
\label{tab:selfrag-matrix}
\begin{tabular}{@{}lllcc@{}}
\toprule
Arm & Models & Isolates & Factoid & Multi-hop \\
\midrule
Monolith-9B        & 1$\times$9B & (reference)           & \CIRCLE & \CIRCLE \\
Decomposed         & 3$\times$4B & topology              & \CIRCLE & \CIRCLE \\
Monolith-4B        & 1$\times$4B & model size            & \CIRCLE & \CIRCLE\texttt{+} \\
Decomposed-Shared  & 1$\times$4B & physical vs.\ logical & \CIRCLE & \LEFTcircle\texttt{+} \\
\midrule
\multicolumn{2}{@{}l}{Engine overlay (HF / vLLM / Ollama)} & inference engine & \Circle & --- \\
\bottomrule
\end{tabular}
\end{table}

\textbf{Metrics.} The full performance set (\S\ref{sec:setup-metrics}),
plus LLM calls per query and resident model memory (defined there); and
the quality set (EM/F1, answered rate, answer parity). Every speed
comparison is \emph{conditioned on} the EM/F1 of the two arms: a speed
result at degraded accuracy is a trade-off, not a free win. This
conditioning is unevaluable on the containment proxy (the arms differ in
verbosity), so the EM/F1 scorer of \S\ref{sec:setup-metrics} is a hard
prerequisite for every E4 table.

\textbf{Collect vs.\ report.} \emph{Collect:} all configs (incl.\ the 4 new
multi-hop controls), both schedules, both DUTs, $R{=}5$, plus the engine
overlay on the single-instance arms. \emph{Report (core):} Monolith-9B
vs.\ Decomposed on \emph{both} tasks and DUTs (the
speed-on-factoid / quality-on-multi-hop trade-off), with Monolith-4B as a
\emph{core} control on both tasks and Decomposed-Shared as a control on
factoid (and multi-hop if collected clean). \emph{Report if clean:} the
Decomposed-Shared multi-hop cell and the engine overlay --- cut first if
space is tight or the MPS ceiling muddies comparability.

% =====================================================================
\subsection{E5 --- MLPerf Inference as a special case (scenario reduction)}
\label{sec:e5-mlperf}

\textbf{Question.} The paper argues single-stream suites miss end-to-end
and collocation effects (Section~\ref{sec:intro}). Rather than compete with
MLPerf, we show it is a \emph{special case}: each MLPerf Inference load
scenario is recoverable as a \sysname{} scheduler configuration. Two of the
four scenario generators exist today (SingleStream = closed-loop,
Server = Poisson); once the saturating-Offline and fixed-interval
MultiStream generators are added (Table~\ref{tab:mlperf-map}), \sysname{}
is a strict superset of the measurement MLPerf performs.

\textbf{Design.} We do \emph{not} claim a certified-throughput parity
number against published LoadGen results on other silicon --- that is
apples-to-oranges and unattributable. Instead the claim is a
\emph{reduction}: MLPerf's four scenarios map onto \sysname{}'s load
generators (Table~\ref{tab:mlperf-map}). We demonstrate the reduction on a
canonical MLPerf vision workload, ResNet-50-v1.5
(\texttt{pipeline\_configs/mlperf/resnet\_inference.yml}), by instantiating
each scenario as a config and confirming \sysname{} reproduces the
scenario's load pattern (verified from the arrival/served-timestamp trace)
and reports the scenario's native metric. \emph{Note:} the committed
resnet config uses \texttt{PoissonLoadScheduler rate=10} --- a
\emph{Server}-like stream, not \emph{Offline}; the Offline cell requires
the saturating all-at-once scheduler. Fixing this mapping is part of the
experiment.

\begin{table}[t]\centering\small
\caption{MLPerf Inference scenarios as \sysname{} configurations. The
reduction, not a throughput comparison, is the claim.}
\label{tab:mlperf-map}
\begin{tabular}{@{}lll@{}}
\toprule
MLPerf scenario & \sysname{} load generator & Native metric \\
\midrule
SingleStream & closed-loop, 1 query in flight (\texttt{OfflineLoadScheduler}$^\dagger$) & p90 latency \\
MultiStream  & fixed $N$-sample query at fixed interval [[FILL: component]] & p99 query latency \\
Server       & Poisson at target QPS (\texttt{PoissonLoadScheduler})     & SLO-bounded QPS \\
Offline      & saturating: all samples enqueued at once [[FILL: component]] & throughput \\
\bottomrule
\end{tabular}
\\[2pt]
\raggedright\footnotesize $^\dagger$Naming hazard: \sysname{}'s
\texttt{OfflineLoadScheduler} is the \emph{closed-loop, one-in-flight}
generator, so it maps to MLPerf \emph{SingleStream} --- \emph{not} MLPerf
\emph{Offline}. The Offline row needs a distinct saturating
(all-at-once) generator and the MultiStream row a fixed-interval
$N$-sample generator; both must exist as \texttt{loadgen/} components or
those rows are claims without an arm ([[FILL: confirm or list to-build]]).
\end{table}

\emph{Optional parity check.} On a single DUT we additionally run the
official MLPerf LoadGen in Offline mode against the same
model/backend and report \sysname{}-within-$X\%$ of the certified
harness on the \emph{same hardware} --- the only parity comparison that
speaks to measurement fidelity ([[FILL: run or explicitly scope out]]).

\textbf{Configurations.}
\texttt{pipeline\_configs/mlperf/resnet\_inference.yml} instantiated once
per scenario by swapping only the \texttt{loadgen} block
(Table~\ref{tab:mlperf-map}); $R{=}5$ per scenario, both DUTs. Three of the
four scenario configs (SingleStream, MultiStream, Offline) are to be
created --- the committed config is the Server-like Poisson cell.
[[FILL: confirm a saturating-Offline and a fixed-interval MultiStream
scheduler exist under \texttt{loadgen/}, else build them.]]

\textbf{Metrics.} Per scenario, the scenario's native metric plus a
trace-level confirmation that the arrival process matches the scenario
definition. If the parity check is run: same-DUT throughput delta vs.\
certified LoadGen.

\textbf{What it shows.} MLPerf Inference is the single-model, isolated
corner of \sysname{}'s configuration space: the \emph{load pattern} of each
scenario is reproducible as a \sysname{} config (confirmed from the arrival
trace), while the end-to-end and collocated regimes (E3, E4, E6) lie
outside MLPerf's expressible space entirely. This is the concrete form of
the paper's ``MLPerf falls short'' thesis.

\textbf{Threats to validity.} The reduction is a claim about the
\emph{load pattern}, verified from the arrival/served-timestamp trace ---
not about bit-exact equivalence with the certified harness, whose
coalescing, tie-break, and warm-up rules differ; recovering a scenario's
certified \emph{number} is exactly what the optional same-DUT LoadGen run
checks. The MultiStream and saturating-Offline generators must exist as
components or the corresponding rows are unbacked. The native-metric column
of Table~\ref{tab:mlperf-map} names MLPerf's certified per-scenario
percentiles (p90 SingleStream, p99 MultiStream); these are the certified
harness's definitions, distinct from \sysname{}'s own pooled-quantile
reporting gate (\S\ref{sec:setup-metrics}, which does not estimate p99 at
our sample sizes) --- computing them at MLPerf's percentiles is part of the
optional parity run, not the reduction itself.

\textbf{Collect vs.\ report.} \emph{Collect:} the four scenario configs on
ResNet-50 both DUTs; optionally the LoadGen parity run. \emph{Report:} the
reduction table with trace-confirmation, and the parity line if collected.
Deliberately small --- a pointed methodological demonstration, not a second
benchmark suite.

% =====================================================================
\subsection{E6 --- Collocation scaling (inter-pipeline interference)}
\label{sec:e6-collocation}

\textbf{Question.} The framework's headline claim is that it can launch and
profile an arbitrary number of concurrent pipelines and attribute
resource use \emph{per process}. E3/E4 collocate stages \emph{within} one
pipeline; here we test the \emph{inter}-pipeline case directly: as an
independent background workload is added to a foreground inference service
on shared hardware, how does foreground quality-of-service degrade, and
can \sysname{} attribute the interference to each pipeline? This is the
capability single-stream, single-workload suites structurally cannot
express.

\textbf{Design.} A single latency-sensitive \emph{foreground} pipeline ---
the EfficientNetV2-S \emph{inference} stream (the inference arm of the
committed \texttt{torchvision\_mixed.yml}), driven by a fixed Poisson rate
--- runs while we vary a \emph{background} load along two axes:
\begin{itemize}
  \item \emph{Concurrency count:} $B\in\{0,1,2,4,\dots\}$ background
    pipelines, each a \emph{separate OS process with its own
    data-loading stage} (the E2 EfficientNetV2-S fine-tune). This
    independence matters for attribution: \texttt{torchvision\_mixed.yml}
    co-locates its inference and training pipelines but they share
    \texttt{dataset\_stage\_id: 0}, so the $B$-generator must give each
    background its own data pipeline, or per-process attribution is
    contaminated by a shared stage.
  \item \emph{Background intensity} (secondary): at fixed $B{=}1$, sweep
    the background batch size / arrival rate to separate ``a co-runner
    exists'' from ``how hard it pushes''.
\end{itemize}
We fix a \emph{single} foreground workload for the reported curve (chosen
once, not ``ResNet-50 or VQA''). The foreground Poisson rate is
\emph{held fixed} across all cells so only the background changes; each
cell runs $R{=}5$ with $\geq\num{500}$ pooled foreground queries so the
p95-vs-$B$ tail points clear the quantile gate of
\S\ref{sec:setup-metrics}. Both DUTs (the 128\,GB GB10 admits far higher
$B$ before saturation than the 16\,GB M2~Pro --- itself a portability data
point).

\textbf{Configurations.} \texttt{torchvision\_mixed.yml} generalised to a
foreground + $B$ \emph{independent} background pipelines, each with its own
dataset stage [[FILL: config generator for the $B$ sweep]]. Per-process
listeners give per-pipeline resource use.

\textbf{Metrics.} Foreground \textbf{p95 and median latency and
throughput} as a function of $B$ (the degradation \emph{curve}, not a
single pair), background throughput, and per-pipeline resource
attribution. The isolation baseline is the $B{=}0$ foreground run at the
same Poisson rate. \emph{Attribution caveat (as in E3).} Per-process
occupancy is \emph{residency}, not utilization: two processes sharing one
GPU are both resident, so occupancy evidences co-residency and bounds
interference but does not causally partition the slowdown. We therefore
report the causal per-pipeline share on GB10 from \texttt{nvidia-smi pmon}
per-process SM activity, and cross-check that the per-process shares
reconcile against the aggregate device counter (a faithfulness check on
the attribution). On the M2~Pro --- which exposes no per-process GPU/ANE
utilization counter --- the degradation \emph{curve} stands but attribution
is explicitly residency-only.

\textbf{What it shows.} A degradation \emph{curve} with \emph{validated}
per-process attribution characterises the collocation-aware-profiling
capability generally (not as one existence-proof cell), and directly
exercises the ``arbitrary number of concurrent pipelines'' claim. It is the
load-bearing evidence for the collocation contribution that
methodology.tex's reviewer note D1 flags as otherwise only asserted.

\textbf{Threats to validity.} The $B$ ceiling is bounded by device memory
and reported, never extrapolated. Attribution is causal only where a
per-process utilization counter exists (GB10); on the M2~Pro it is
residency-only. Backgrounds must be independent processes with independent
dataset stages, or attribution is contaminated. The foreground rate is
held fixed so the $x$-axis is purely $B$.

\textbf{Collect vs.\ report.} \emph{Collect:} the full $B$ sweep on both
DUTs plus the intensity sub-sweep. \emph{Report:} the foreground p95-vs-$B$
degradation curve with per-process occupancy on both DUTs; the intensity
sub-sweep if space allows. \emph{Scope honesty:} $B$ ranges are bounded by
device memory --- we report the ceiling reached on each DUT rather than
implying unbounded scaling.

% =====================================================================
\subsection{E7 --- Capacity portability sweep (no-code scalability)}
\label{sec:e7-sweep}

\textbf{Question.} The paper claims one can stress hardware of any capacity
by editing a declarative \texttt{model} field --- no code change. E4's
9B-vs-4B is two points and confounded with topology. Here we test the
scalability claim cleanly: does a \emph{single} pipeline, reconfigured only
by its model field, sweep from comfortably-resident to
capacity-saturating on each DUT, and where does each device's ceiling fall?

\textbf{Design.} We do \emph{not} run a separate pipeline for this: E7 is
the Self-RAG \emph{monolith} arm of E4 read along the size axis. The two
monolith configs (\texttt{factoid\_monolith\_*} at 9B and
\texttt{factoid\_monolith\_4b\_*} at 4B) are \emph{byte-identical except the
\texttt{model} field} --- same stages, same retrieval collection, same
loadgen --- so they already constitute a config-only, same-pipeline
2-point size comparison at fixed topology (this is exactly the ``size vs.\
topology'' control of E4). E7 extends that ladder with additional rungs by
swapping \emph{only} the \texttt{model} field. Quantization is matched
\emph{within} each DUT (M2~Pro all-4-bit MLX; GB10 all-BF16) and the ladder
is read \emph{within} a device; context length and in-flight concurrency
are held fixed so the size axis is not confounded. To make the
\emph{no-code} claim concrete we show the literal config diff between the
smallest and largest rung --- it is the \texttt{model} field and nothing
else (verified for the committed 4B/9B pair). On each DUT the ceiling is
the largest size that fits with \emph{zero} OOM across all $R{=}5$ runs.

\emph{The ladder is concrete and genuinely config-only.} The Qwen3.5
family ships OptiQ-4bit MLX builds across a full dense range that
\emph{share one chat template and architecture}, so a rung change is only
the \texttt{model} field --- no tokenizer, template, or recipe change. On
the M2~Pro (4-bit) the dense ladder is 0.8B (\SI{0.5}{\giga\byte}), 2B
(\SI{1.4}{\giga\byte}), 4B (\SI{2.8}{\giga\byte}), 9B
(\SI{5.6}{\giga\byte}), 27B (\SI{15.7}{\giga\byte}): the first four fit and
\textbf{27B is the ceiling rung --- its \SI{15.7}{\giga\byte} of weights
alone exceed the usable \SI{16}{\giga\byte} budget, so it OOMs} while GB10
runs it with headroom. The 4B/9B rungs come free from E4's size control;
0.8B, 2B, and 27B are one-line additions. We keep the sparse
\texttt{35B-A3B} MoE \emph{off} the dense ladder --- mixing an MoE into a
dense parameter-size axis would confound size with sparsity --- and report
it, if at all, as a separate point. The free 4B/9B pair already
substantiates the \emph{no-code} half; the 27B rung supplies the capacity
\emph{ceiling}. Because the monolith carries the retry/rewrite loop,
per-query \emph{latency} is not a clean function of size on this pipeline:
the clean Self-RAG size axes are \textbf{resident memory} and
\textbf{answer quality} vs.\ size (the latter corroborated by the family's
published capability scores, which rise $36.0\!\to\!79.1$ from 0.8B to
27B). A clean throughput/latency-vs-size \emph{curve}, if wanted, is better
taken from the non-agentic VQA or a plain-generation pipeline.

\textbf{Configurations.} The E4 monolith configs at 4B and 9B (already
collected), plus three added \texttt{model}-field rungs on M2~Pro ---
\texttt{Qwen3.5-\{0.8B,2B\}-OptiQ-4bit} (fit) and
\texttt{Qwen3.5-27B-OptiQ-4bit} (the \SI{15.7}{\giga\byte} ceiling rung,
expected to OOM) --- and their BF16 counterparts on GB10; $R{=}5$ per size
that fits.

\textbf{Metrics.} Per size: resident model memory, per-query latency,
throughput; and per DUT the ceiling (largest size with zero OOM across
$R{=}5$ runs at fixed context/concurrency) --- the 16\,GB M2~Pro ceiling
vs.\ the 128\,GB GB10 headroom.

\textbf{What it shows.} A config-only size ladder that saturates the small
device while the large one still has headroom is the most direct evidence
for the ``mid-range gap / portability'' argument (Section~\ref{sec:background})
and the ``no-code scalability'' contribution --- currently asserted from the
8$\times$/1.4$\times$ capacity/bandwidth framing but never exercised to a
limit.

\textbf{Threats to validity.} ``Fits'' is a (model, quant, context,
concurrency) property, not a device constant; we pin quant (within DUT),
context, and concurrency so the reported ceiling is reproducible and the
size axis is unconfounded. The ``no-code'' claim is substantiated by the
literal config diff \emph{and} by the family sharing one chat
template/architecture across the dense ladder, so a rung change is
genuinely only the \texttt{model} field. The sparse MoE is kept off the
dense size axis to avoid confounding parameter count with sparsity.

\textbf{Collect vs.\ report.} \emph{Collect:} the 4B/9B points come free
from E4's monolith arm; only the \emph{added rungs} (and the
ceiling-reaching large rung, if a 4-bit build exists) are new collection.
\emph{Report:} the resident-memory- and quality-vs-size points with each
device's ceiling marked. \emph{Reviewer-driven addition:} E7 exists to
back the scalability claim; if the Intro is instead softened to drop
``any-size hardware,'' E7 can fold into E4's discussion rather than stand
alone.

% =====================================================================
\subsection{Summary: what we collect vs.\ what we report}
\label{sec:setup-summary}

The collection plan is broader than the paper. We \emph{collect}: the
re-collected NoOp matrix (E1); baseline+tracing arms on both DUTs (E2); the
VQA $2\times2$ plus DRAM-bandwidth trace (E3); the full Self-RAG matrix
with the size, logical-decomposition, and engine controls --- including the
four multi-hop control configs to be created --- on both DUTs (E4); the
ResNet-50 scenario configs (E5); the collocation $B$-sweep (E6); and the
model-size ladder (E7). We \emph{report}, as the paper's spine: two
instrument-fidelity bounds (E1, E2) that license attributing later
overhead to workloads; two case studies whose primary claims are the
contention-bound heterogeneity collapse (E3) and the config-only
monolith-vs-decomposition trade-off with its size and physical-vs-logical
controls (E4); and three capability results --- the MLPerf reduction (E5),
the collocation degradation curve (E6, the load-bearing collocation
evidence), and the no-code capacity sweep (E7). Secondary arms (the
Self-RAG engine overlay, the Decomposed-Shared multi-hop cell, the E6
intensity sub-sweep) are reported only where the signal is clean.

Collection dependencies gate the final tables and are tracked as blockers.
\emph{Hard (a claim is unreportable without it):} the \textbf{EM/F1 (and
VQA-accuracy) scorers}, without which no accuracy-conditioned comparison
(E3, E4) can be reported; the \textbf{GB10 halves} of E2/E4/E6/E7, without
which every cross-device direction claim is one-sided; and the
\textbf{re-collection of E1} on an idle M2~Pro, since the committed matrix
violates the single-quiet-machine premise. \emph{Build work (new framework
code / config generators the capability tier needs):} the two missing
\textbf{\texttt{loadgen/} scheduler components} --- a saturating-Offline and
a fixed-interval MultiStream generator --- which back 2 of the 4 rows of the
E5 reduction table; the \textbf{E6 $B$-sweep config generator} that spawns
independent background processes each with its own dataset stage; and the
\textbf{E7 model-size ladder configs}. \emph{Smaller items:} the
\textbf{four multi-hop Self-RAG control configs} (E4), the
\textbf{E3 DRAM-bandwidth read} (without which E3 drops from
``contention-bound'' to ``consistent with contention''), and the
\textbf{synchronous tracing cost} from an orchestrated run
(E1/E2 overhead-in-context). Finally, this draft delivers
the collocation (E6) and MLPerf-contrast (E5) experiments that
\texttt{methodology.tex} currently defers to future work; that file's
Experiments subsection and the Intro's scalability wording must be
reconciled when this draft replaces it.
